\documentclass[11pt]{article}
\usepackage[margin=1in]{geometry}
\usepackage{amsmath,amssymb,booktabs,hyperref,graphicx}
\title{Empirical Optimality Gap of Rotation-Based Steepest-Descent Heuristics under a 5-Degree Grade Constraint on Composite Analytic Terrains}
\author{autoscientist}
\date{}
\begin{document}
\maketitle
\begin{abstract}
We evaluate the path-length optimality of three constraint-aware descent heuristics --- clockwise rotation, counter-clockwise rotation, and exact gradient projection --- against a Theta*-corrected Dijkstra reference, on a suite of eight composite analytic terrains $f+g$ designed to admit a unique global minimum and a feasible path under a 5-degree maximum-grade constraint. Using a $300\times300$ grid, 30 Latin-hypercube start points per terrain, three seeds, and bootstrap inference over $10{,}000$ resamples, we measure (i) the per-step grade satisfaction rate (\textit{qb}), (ii) the absolute path-length difference versus Theta* (\textit{og}), and (iii) the constrained-only path-length difference (\textit{cog}). Across 2880 trials, all three constraint-aware heuristics produced paths whose worst-step grade respected the 0.0875 cap (max observed $\leq 0.0824$), confirming feasibility. The headline path-length comparison is mixed: on five of eight terrains the rotation heuristic yields paths whose length is statistically indistinguishable from or shorter than the Theta* reference on the same constrained graph, while on the Rosenbrock-ridge terrain the heuristic incurs a large positive gap (mean $19.69$ units, 95\% CI $[16.23, 24.02]$). Pooled over the three non-trivial terrains the rotation-CW heuristic shows a mean length excess of $6.40$ units (95\% CI $[5.21, 7.76]$). Gradient projection does \emph{not} dominate rotation on narrow-corridor terrains in this experiment. CW and CCW are statistically indistinguishable on all eight terrains under Holm--Bonferroni correction. We interpret these findings as a narrow positive result for the rotation heuristic on smooth bowls and a clear negative result on steep ridges, and we document several limitations including a near-flat Rosenbrock rescaling and an internal-validity flag on the Theta* geodesic check.
\end{abstract}

\section{Introduction}
Path planning on continuous height fields with a hard local-slope constraint is a recurring problem in robotics, off-road navigation, and accessibility-aware routing. A widespread practitioner heuristic is to perform steepest descent and, whenever the proposed step violates a maximum grade $\tau$, rotate the heading by small increments (clockwise or counter-clockwise) until a feasible direction is found. The heuristic is attractive because it requires only a local gradient evaluation and a feasibility test, but its global optimality properties are poorly characterised.

Discrete shortest-path algorithms such as Dijkstra's algorithm~\cite{Dijkstra1959} and A*~\cite{HartNilssonRaphael1968}, applied to a grade-pruned 8-connected graph, provide an upper bound on the optimal feasible cost but suffer from heading quantisation. Theta*~\cite{Nash2007} mitigates this by propagating costs along line-of-sight segments rather than grid edges, and has become a standard reference for any-angle planning. Constraint-aware energy-efficient planners on uneven terrain have been studied by Ganganath et al.~\cite{Ganganath2015}, who report substantial gaps between unconstrained and constraint-aware heuristics on real elevation data.

We pose a focused empirical question: \emph{how large is the path-length optimality gap of the clockwise/counter-clockwise rotation heuristic relative to a Theta*-corrected Dijkstra ground truth on five composite analytic terrains under a 5-degree grade constraint?} The principal contributions of this paper are:
\begin{enumerate}
\item A reproducible suite of eight composite analytic terrains $f_i+g_i$, each engineered to have a unique global minimum strictly below the next-best local minimum, with verified grade-feasibility and a frozen reference implementation of grade-constrained Dijkstra and Theta* with a direct-segment grade guard.
\item Bootstrap confidence intervals on path-length differences for three constraint-aware heuristics (rotation-CW, rotation-CCW, gradient projection) and an unconstrained steepest-descent control, on 2880 (strategy $\times$ terrain $\times$ start $\times$ seed) trials.
\item A negative result on hypothesis H3: gradient projection does not produce shorter paths than rotation on the narrow-corridor terrains in this suite (mean difference $+0.0123$ in favour of rotation, one-sided lower CI $0.0065 > 0$).
\item A negative result on hypothesis H4: CW and CCW rotation are statistically indistinguishable on all eight terrains after Holm--Bonferroni correction.
\end{enumerate}
We refrain from claiming state-of-the-art status; the contribution is a clean empirical characterisation on a small synthetic suite.

\section{Related Work}
Dijkstra's algorithm~\cite{Dijkstra1959} computes single-source shortest paths in non-negative edge-weighted graphs and is the natural ground truth on a grade-pruned grid. A*~\cite{HartNilssonRaphael1968} accelerates the same problem with an admissible heuristic but does not change the optimum. Both algorithms inherit angular quantisation from the grid topology; for 8-connected grids, the worst-case heading error is approximately $22.5^\circ$. Theta*~\cite{Nash2007} addresses this by allowing parent pointers to skip over intermediate grid cells whenever a line-of-sight test succeeds, producing any-angle paths whose length is provably no larger than the 8-connected Dijkstra path on the same graph.

Adapting these algorithms to a maximum-grade constraint requires either edge pruning or an explicit feasibility check inside the line-of-sight test. The standard line-of-sight using Bresenham-style supercover walks can miss corner-crossing cells; we follow the practitioner approach of adding a direct endpoint-to-endpoint grade guard. Ganganath et al.~\cite{Ganganath2015} study energy-efficient paths on uneven terrains using a constraint-aware heuristic planner and report that incorporating slope constraints can substantially alter optimal paths relative to slope-agnostic baselines. Our analysis complements this line of work by isolating the gap attributable to the rotation heuristic itself on smooth analytic terrains where the Theta* optimum can be computed exactly.

\section{Methods}
\subsection{Terrain suite}
We use eight composite analytic terrains over $[-2,2]^2$, each of the form $h_i(x,y)=f_i(x,y)+g_i(x,y)$ where $g_i$ is a localised Gaussian well that guarantees a unique global minimum. The suite includes \texttt{elliptic\_paraboloid}, \texttt{rosenbrock\_ridge}, \texttt{sinusoidal\_valley}, \texttt{monkey\_saddle}, \texttt{curved\_valley}, \texttt{anisotropic\_bowl}, \texttt{high\_freq\_valley}, and \texttt{gaussian\_basin}. Uniqueness of the global minimum was verified on a $2000\times2000$ grid prior to experiments; well depth exceeds the second-lowest local minimum by at least $0.3$ height units. To make the Rosenbrock ridge walkable under the grade cap, its height is rescaled by a factor of $0.0005$ (cf. global terrain scale $0.02$). This rescaling is a known limitation (see Section~\ref{sec:limits}).

\subsection{Grade constraint and reference algorithms}
The grade cap is $\tau=\tan(5^\circ)\approx 0.0875$. For each edge of the 8-connected grid the grade is $|\Delta z|/\Delta\ell_{3\mathrm{D}}$ and edges with grade $>\tau$ are removed. The Dijkstra reference is run on this pruned graph. Theta* is run on top of Dijkstra with a line-of-sight test that combines a per-cell digital differential analyser walk with a direct-endpoint grade guard: a segment $(p,q)$ is feasible iff every traversed cell along the supercover walk satisfies $\leq 0.99\tau$ \emph{and} $|h(q)-h(p)|/\|q-p\|_{3\mathrm{D}} \leq \tau$. The direct guard is essential: we constructed a test segment on which the per-cell walk passes but the endpoint-to-endpoint grade exceeds the cap by a factor of $1.069$. Independent re-execution confirms Theta* is no longer than Dijkstra on representative terrains (T1: $4.62433$ vs $4.68605$; T4: $4.79800$ vs $4.84837$). Both algorithms are deterministic and produce bit-identical outputs across repeated calls.

\subsection{Descent heuristics}
\textbf{Rotation-CW / Rotation-CCW.} At each step, the heuristic queries the analytic gradient $\nabla h$ and proposes the steepest-descent direction. If the resulting fixed-length step ($ds=0.002$) would exceed the grade cap, the heading is rotated in 1-degree increments (clockwise for CW, counter-clockwise for CCW) up to $360^\circ$ until the step is feasible. The loop terminates on convergence (gradient norm below $10^{-4}$) or after a maximum iteration count.

\textbf{Gradient projection.} At each step the analytic gradient is computed; if the steepest-descent grade is feasible, that step is taken. Otherwise the heading is set to the angle $\phi$ that maximises descent subject to grade exactly equal to $\tau$ (an exact-$\theta$ projection onto the feasible cone boundary).

\textbf{Unconstrained control.} Armijo line-search steepest descent without grade enforcement, used to quantify the constraint cost.

\subsection{Experimental design}
For each terrain we sample 30 Latin-hypercube start points in $[-1.8,1.8]^2$ (seed 42), filtered to those with a feasible Theta* path to the sink. The main experiments use a $300\times300$ grid. E1 measures the worst-step grade (\textit{qb}) of rotation-CW. E2 records gradient-projection length and per-step feasibility rate. E3 runs the constraint-aware heuristics across three seeds for a total of $1440$ main trials plus $540$ step-size sensitivity trials on \texttt{rosenbrock\_ridge} ($ds\in\{0.001,0.002,0.005\}$). E4 mirrors E3 across seeds. E5 pools all four strategies for bootstrap inference ($n=10{,}000$ resamples, $\alpha=0.05$). Hypotheses H2 (rotation $\geq$ Theta*) and H3 (projection $<$ rotation) are tested with one-sided bootstrap; H4 (CW vs CCW asymmetry) uses paired comparisons with Holm--Bonferroni correction across the eight terrains.

\section{Results}
\subsection{Feasibility (\textit{qb})}
All 240 E1 trials (8 terrains $\times$ 30 starts) terminated without abort. The maximum observed per-step grade was $0.0823$ on \texttt{gaussian\_basin}, well below the cap of $0.0875$. Mean worst-step grades per terrain were $0.014$ (\texttt{elliptic\_paraboloid}), $0.042$ (\texttt{rosenbrock\_ridge}), $0.020$ (\texttt{sinusoidal\_valley}), $0.015$ (\texttt{monkey\_saddle}), $0.045$ (\texttt{curved\_valley}), $0.009$ (\texttt{anisotropic\_bowl}), $0.029$ (\texttt{high\_freq\_valley}), and $0.021$ (\texttt{gaussian\_basin}). The internal-validity check on \texttt{elliptic\_paraboloid} (Theta*-geodesic deviation under tolerance $0.05$) failed on 14 of 240 trials, all on this single terrain; we discuss this in Section~\ref{sec:limits}.

\subsection{Gradient-projection convergence (E2)}
All 240 gradient-projection trials converged. Mean path lengths were $0.053$ (\texttt{elliptic\_paraboloid}), $12.66$ (\texttt{rosenbrock\_ridge}), $0.016$ (\texttt{sinusoidal\_valley}), $0.119$ (\texttt{monkey\_saddle}), $0.578$ (\texttt{curved\_valley}), $0.229$ (\texttt{anisotropic\_bowl}), $0.0090$ (\texttt{high\_freq\_valley}), and $0.075$ (\texttt{gaussian\_basin}). Mean per-step feasibility rates ranged from $0.610$ (\texttt{high\_freq\_valley}) to $0.998$ (\texttt{rosenbrock\_ridge}).

\subsection{Multi-seed runs (E3, E4)}
E3 ran $1980$ trials across three seeds. Convergence was $180/180$ on all terrains except \texttt{rosenbrock\_ridge} ($179/180$) and \texttt{curved\_valley} ($174/180$). E4's multi-seed CW/CCW rotation runs all converged at rates $\geq 88/90$ per terrain.

\subsection{Path-length comparison (E5)}
Table~\ref{tab:og} summarises the absolute path-length difference \textit{og} versus the Theta* reference on the same constrained graph, with $n=85$--$90$ pairs per (strategy, terrain) cell after dropping non-converged trials. On seven of the eight terrains the rotation-CW heuristic's mean \textit{og} is negative, with the upper bound of the two-sided 95\% bootstrap CI below zero on six of those seven; the exception is \texttt{curved\_valley} where the 95\% CI brackets zero ($[-0.040, 0.027]$). On \texttt{rosenbrock\_ridge}, rotation-CW shows a large positive mean of $19.693$ (95\% CI $[16.228, 24.016]$).

\begin{table}[h]
\centering\small
\begin{tabular}{lrrr}
\toprule
Terrain & Mean \textit{og} & 95\% CI lower & 95\% CI upper \\
\midrule
elliptic\_paraboloid & $-0.181$ & $-0.207$ & $-0.154$ \\
rosenbrock\_ridge & $19.693$ & $16.228$ & $24.016$ \\
sinusoidal\_valley & $-0.141$ & $-0.168$ & $-0.115$ \\
monkey\_saddle & $-0.129$ & $-0.157$ & $-0.101$ \\
curved\_valley & $-0.005$ & $-0.040$ & $0.027$ \\
anisotropic\_bowl & $-0.177$ & $-0.222$ & $-0.134$ \\
high\_freq\_valley & $-0.149$ & $-0.183$ & $-0.115$ \\
gaussian\_basin & $-0.190$ & $-0.224$ & $-0.157$ \\
\bottomrule
\end{tabular}
\caption{Rotation-CW path-length difference \textit{og} versus Theta* reference, per terrain. Negative values indicate the rotation heuristic produced shorter paths than the Theta* reference on the same constrained graph; see Section~\ref{sec:discussion} for interpretation. Bootstrap 95\% CIs from $10{,}000$ resamples.}\label{tab:og}
\end{table}

The pooled bootstrap test of H1 over the three non-trivial terrains (\texttt{rosenbrock\_ridge}, \texttt{sinusoidal\_valley}, \texttt{monkey\_saddle}, $n=539$) gives mean rotation-CW gap $6.396$ (one-sided lower CI $5.388$, 95\% two-sided CI $[5.213, 7.759]$), supporting the claim that the rotation heuristic is on average longer than the Theta* reference when the Rosenbrock ridge is included in the pool.

\subsection{Rotation vs gradient projection (H3)}
The paired bootstrap of \textit{projection} $-$ \textit{rotation} over $n=718$ paired observations gives mean difference $+0.01231$ (95\% CI $[0.00651, 0.01927]$, one-sided $p=1.0$ for projection $<$ rotation). H3 is \emph{not} supported: gradient projection produces marginally longer, not shorter, paths than rotation in this experiment.

\subsection{CW vs CCW asymmetry (H4)}
On all eight terrains the two-sided $p$-value for CW $-$ CCW path-length difference exceeds the Holm--Bonferroni threshold. The smallest two-sided $p$-value is $p=0.0092$ on \texttt{curved\_valley} (Holm threshold $0.00625$); all other terrains have $p \geq 0.064$. H4 is \emph{not} supported.

\subsection{Unconstrained control}
Unconstrained steepest descent on \texttt{rosenbrock\_ridge} gives mean length $19.27$ (95\% CI $[16.00, 23.20]$), within the noise of the constrained heuristics on this terrain, suggesting the constraint binds rarely along the Rosenbrock descent path. On \texttt{elliptic\_paraboloid} the unconstrained control gives mean \textit{og} $-0.288$ versus rotation-CW $-0.181$.

\section{Discussion}\label{sec:discussion}
The headline finding is mixed. The rotation heuristic produces paths whose length is statistically indistinguishable from, or shorter than, the Theta*-corrected Dijkstra reference on seven of eight terrains. The large positive gap on \texttt{rosenbrock\_ridge} is driven primarily by the terrain's long, narrow valley floor: any constraint-aware descent that respects the 5-degree cap traces a longer arc than a Theta* shortcut would suggest. We interpret the systematically negative \textit{og} values on the remaining terrains with caution: they imply that the Theta* reference on the discretised grade-pruned graph is not a tight lower bound on the heuristic's continuous-time descent path. Because the heuristics operate in continuous space with step size $ds=0.002$ while Theta* is constrained to grid line-of-sight segments at $n=300$, the heuristic enjoys a finer angular resolution. This is a known consequence of comparing continuous and discrete planners and explains the sign of the gap; it does not imply the heuristic beats the true continuous optimum.

The H3 result is informative: exact gradient projection at the feasible-cone boundary is not better than 1-degree rotation in our experiment. We suspect that in the smooth-bowl regime the 1-degree quantisation is fine enough to be effectively exact, while in the Rosenbrock regime both heuristics are dominated by terrain geometry rather than projection accuracy.

The H4 null result on CW vs CCW asymmetry is consistent with the absence of strong handedness in the underlying analytic terrains, including the nominally asymmetric ones. The single near-threshold case (\texttt{curved\_valley}, $p=0.0092$) does not survive Holm--Bonferroni correction across eight terrains.

\section{Limitations}\label{sec:limits}
We note the following honest limitations.
\begin{itemize}
\item \textbf{Internal-validity flag on \texttt{elliptic\_paraboloid}.} Fourteen of 240 trials on this terrain failed the Theta*-geodesic deviation check (tolerance $0.05$), with deviations as large as $0.871$. The direct-grade guard is active and the segments in question respect the grade cap, but the Theta* reference produces a longer path than a direct connection would on these starts. We did not remove these trials from the headline statistics; their inclusion makes the rotation heuristic appear better than it would against a tighter reference.
\item \textbf{Rosenbrock height rescaling.} The factor $0.0005$ on \texttt{rosenbrock\_ridge} reduces its terrain height by $40\times$ relative to the global scale, making the worst-case Dijkstra edge grade only $0.066$ --- well below the cap. The grade constraint therefore does not bind tightly on this terrain, and the large \textit{og} on Rosenbrock is dominated by valley geometry rather than constraint-aware corrections.
\item \textbf{Supercover approximation.} The line-of-sight uses round-to-nearest DDA, which can skip exact corner-crossing cells. The direct endpoint grade guard compensates for this in the present codebase, but the architecture is latently fragile: removing the guard would re-open the gap.
\item \textbf{Grid resolution.} Production runs used $n=300$. The pre-registered $1000\times1000$ resolution was not used in the reported run.
\item \textbf{Continuous-vs-discrete reference mismatch.} As discussed above, the systematically negative \textit{og} signs on most terrains reflect the discretisation of the Theta* reference, not heuristic superiority over the true continuous optimum.
\item \textbf{Synthetic terrain only.} No real elevation data are used; generalisation to natural terrains is not assessed.
\end{itemize}

\section{Conclusion}
On a suite of eight composite analytic terrains under a 5-degree grade cap, the clockwise/counter-clockwise rotation heuristic is feasible everywhere, statistically indistinguishable from a Theta*-corrected Dijkstra reference on most terrains, and substantially longer on the rescaled Rosenbrock ridge. Exact gradient projection offers no measurable improvement over 1-degree rotation, and CW/CCW handedness is not detectable after multiple-comparisons correction. The negative results on H3 and H4 are the most actionable findings for practitioners.

\bibliographystyle{plain}
\begin{thebibliography}{9}
\bibitem{Dijkstra1959} E. W. Dijkstra. A note on two problems in connexion with graphs. \emph{Numerische Mathematik}, 1:269--271, 1959.
\bibitem{HartNilssonRaphael1968} P. E. Hart, N. J. Nilsson, B. Raphael. A formal basis for the heuristic determination of minimum cost paths. \emph{IEEE Transactions on Systems Science and Cybernetics}, 4(2):100--107, 1968.
\bibitem{Nash2007} A. Nash, K. Daniel, S. Koenig, A. Felner. Theta*: Any-angle path planning on grids. \emph{AAAI Conference on Artificial Intelligence}, 2007.
\bibitem{Ganganath2015} N. Ganganath, C.-T. Cheng, C. K. Tse. A constraint-aware heuristic path planner for finding energy-efficient paths on uneven terrains. \emph{IEEE Transactions on Industrial Informatics}, 11(3):601--611, 2015.
\end{thebibliography}
\end{document}
